\section{Case Study: Coral Triangle Reef Restoration}

\label{sec:case_study}
% ══════════════════════════════════════════════════════════════════════════════

We present a detailed case study for a BIB financing coral reef
restoration across five sites in the Coral Triangle.

\subsection{Project Parameters}

\begin{itemize}
  \item \textbf{Target}: Increase hard coral cover from 12\% to 25\%
    across 500\,ha over 10 years.
  \item \textbf{Total capital}: \$8.5M (7{,}500 BIB tokens at \$1{,}133).
  \item \textbf{Outcome payer}: National governments (40\%), tourism
    consortium (35\%), Blue Carbon initiative (25\%).
  \item \textbf{Verification}: Quarterly photoquadrat surveys,
    continuous acoustic monitoring, annual satellite turbidity.
\end{itemize}

\subsection{Projected Returns}

Monte Carlo simulation with coral growth models yields:
\begin{itemize}
  \item Expected IRR: 7.8\% (median), 9.4\% (mean).
  \item Full return probability: 68\%.
  \item Partial return probability ($>$50\% of principal): 89\%.
  \item Total loss probability: 4\%.
  \item Ecosystem service value: \$34M over 20 years (4:1 ratio).
\end{itemize}

\subsection{Key Lessons}

Multi-stakeholder outcome payer structures reduce single-payer default
risk. Continuous acoustic monitoring enables early trajectory deviation
warning. Ecosystem service co-benefits (fisheries, tourism, coastal
protection) provide multiple revenue streams for outcome payments.


% ══════════════════════════════════════════════════════════════════════════════
